\section{Higher symplectic cores over the field of five elements}
\label{sec:higher-symplectic-five}

\paragraph{Status and scope.}
This complete new elementary proof has passed independent whole-proof
reviews by /root and /root/imprimitive. Both also read the full author
checker; the peer separately reconstructed the exact finite arithmetic.
Its character-theoretic construction is the explicit
Heisenberg matrix algebra in D-HOS-01, extended here to automorphisms
with fixed vectors and to defining-characteristic prime order. No
unipotent character bound or full character table is used. The final
source obstruction allows mixed ranks, arbitrary compatible outer
automorphisms, scalar twists, and permutations of tensor factors.
It uses the imaginary character-field inclusion from the original
induced representation, not equality of the original and inducing fields.

For a finite-order complex operator $T$ on a space of dimension $a$, write
\[
 \nu(T)=a-\max_\lambda\dim\ker(T-\lambda I).
                                                               \tag{HOS5.1}
\]
This quantity is unchanged by scalar multiplication, is no larger than
ordinary moving rank, and satisfies $\nu(T^j)\le\nu(T)$.

\begin{lemma}[D-HOS-07: opposite line radicals]
\label{D-HOS-07}
Let $q$ be odd, $n\ge2$, and let $V$ be a nontrivial irreducible
ordinary representation of $G=\operatorname{Sp}_{2n}(q)$ of degree $a$.
If $Q$ is the unipotent radical of a line parabolic, then
\[
                         \dim V^Q\le a/2.                     \tag{HOS5.2}
\]
\end{lemma}

\begin{proof}
Let the natural space be $E=\langle e,f\rangle\perp W$, with
$B(e,f)=1$. For $w\in W$ and $z\in\mathbf F_q$ put
\[
 u_e(w,z)e=e,\quad
 u_e(w,z)x=x+B(w,x)e\ (x\in W),\quad
 u_e(w,z)f=f+w+2ze.                                          \tag{HOS5.3}
\]
Direct substitution in the alternating form shows these are symplectic.
Their multiplication is
\[
 u_e(w,z)u_e(w',z')
 =u_e(w+w',z+z'+B(w,w')/2).
\]
They form $Q_e=Q$. Define $Q_f$ by the same formula with the paired
vectors $(f,-e)$ and the same $W$. The symplectic map
$e\mapsto f$, $f\mapsto-e$, $W\mapsto W$ conjugates $Q_e$ to $Q_f$.
We prove directly that these two groups generate $G$.

For $v\in E$ let
\[
                       t_v(c)x=x+cB(v,x)v.                    \tag{HOS5.4}
\]
The centers of $Q_e,Q_f$ contain $t_e(c),t_f(c)$ for every $c$.
Conjugating $t_f(c)$ by (HOS5.3) gives every transvection in a
direction $f+w+de$, with $w\in W,d\in\mathbf F_q$. The opposite
calculation gives every direction $e+w+df$. For $0\ne w\in W$,
choose $w_0\in W$ with $B(w_0,w)=-1$. Then
\[
                         u_f(w_0,0)(f+w)=w.
\]
Consequently transvections in direction $w$ also belong to
$\langle Q_e,Q_f\rangle$. Scaling a direction scales the parameter
by its square, and the parameter already ranges over the whole field.
Every nonzero vector is either in $W$ or a scalar multiple of one of
the displayed directions. Thus this generated group contains all
symplectic transvections.

For completeness these transvections generate $G$ by an elementary
induction. If $B(x,y)\ne0$, a suitable transvection in direction
$y-x$ sends $x$ to $y$. If $B(x,y)=0$, choose a vector outside the
two orthogonal hyperplanes and use two such steps; two proper
hyperplanes cannot cover a vector space over an odd field. Hence
one can first make a given symplectic map fix $e$. Its image of $f$
then has the form $f+w+de$ and can be moved to $f$ by an element of
$Q_e$. These radical elements are themselves products of
transvections: since $B(e,w)=0$, direct multiplication gives
\[
 u_e(w,0)=t_{w+e}(1)t_w(-1)t_e(-1),\qquad
 u_e(w,z)=t_e(2z)u_e(w,0).                                   \tag{HOS5.5}
\]
After fixing both $e$ and $f$, continue on $W$. This includes the
two-dimensional base case, where the remaining orthogonal complement
is zero. It proves $\langle Q_e,Q_f\rangle=G$.

The two fixed subspaces in $V$ have equal dimension. Their intersection
is $V^G=0$, because $V$ is nontrivial and irreducible. Their dimensions
therefore sum to at most $a$, proving (HOS5.2).
\end{proof}

\begin{lemma}[D-HOS-08: carriers with a fixed subspace]
\label{D-HOS-08}
In the notation of D-HOS-07, suppose a finite-order operator $T$ on
$V$ acts on the full line radical by
\[
             T\rho(u_e(w,z))T^{-1}=\rho(u_e(gw,z)),
                                                               \tag{HOS5.6}
\]
where $g\in\operatorname{Sp}(W)$ has prime order $\ell$. Put
$r=\operatorname{rank}(g-I)>0$. The prime $\ell$ may equal the
characteristic. Then
\[
 \nu(T)\ge (1-1/\ell)(1-q^{-r/2})(a-\dim V^{Q_e})
          \ge\tfrac12(1-1/\ell)(1-q^{-r/2})a.                \tag{HOS5.7}
\]
\end{lemma}

\begin{proof}
We retain the actual group law in (HOS5.3), so no choice of projective
intertwiner changes the specified automorphism. For a nontrivial
central additive character $\psi$, D-HOS-01 constructs the unique
irreducible carrier $U_\psi$ of dimension $N=q^{n-1}$ and its
operators $L_w$. These satisfy
\[
 L_wL_{w'}=\psi(B(w,w')/2)L_{w+w'},\qquad
 \{L_w:w\in W\}\text{ is a basis of }\operatorname{End}(U_\psi).
                                                               \tag{HOS5.8}
\]
The proof there is the explicit translation-and-multiplication model
and the central-character group-algebra dimension calculation; it
does not use the order of $g$ or a fixed-point-free hypothesis.
Uniqueness gives an exact intertwiner
$\Omega L_w\Omega^{-1}=L_{gw}$. Its $\ell$th power is scalar;
rescale it so that $\Omega^\ell=I$.

Every $g^j$, $1\le j<\ell$, has the same fixed space as $g$.
Indeed $g$ is a power of $g^j$ in the cyclic group of prime order.
Conjugation by $\Omega^j$ permutes the basis in (HOS5.8), whence
\[
 |\operatorname{Tr}\Omega^j|^2
 =\operatorname{Tr}(\operatorname{Ad}\Omega^j)
 =q^{2n-2-r},\qquad
 |\operatorname{Tr}\Omega^j|/N=q^{-r/2}.                     \tag{HOS5.9}
\]
The first identity is valid because $\Omega$ has finite order.
Fourier inversion bounds any one eigenspace of $\Omega$ by
\[
              \frac{1+(\ell-1)q^{-r/2}}{\ell}\,N.
                                                               \tag{HOS5.10}
\]
On the full $\psi$-isotypic summand, the operator is
$\Omega\otimes A_\psi$ by the full matrix-algebra commutant.
The multiplicity operator is finite order after the above
normalization, since a common multiple of the orders of $T$ and
$\Omega$ kills it. Its eigenvalue multiplicities convolve with
those of $\Omega$, so (HOS5.10) remains an upper bound on the
normalized multiplicity. No scalar-power hypothesis on $T$ over
the whole of $V$ is being made.

The remaining nontrivial irreducibles of $Q_e$ are its nonzero
linear characters, parametrized by the nonzero vectors in the
dual of $W$ using the finite-field trace pairing. Their
multiplicities in $V$ are all equal: the subgroup
$\operatorname{Sp}(W)\le G$ fixing $e,f$ normalizes $Q_e$ and is
transitive on these nonzero dual vectors. Of the $q^{2n-2}-1$
weights, exactly $q^{2n-2-r}-1$ are fixed by $g$. Their proportion
is less than $q^{-r}$. Every other orbit has length $\ell$;
on such a cyclic orbit each eigenvalue has multiplicity at most
$1/\ell$ of the total orbit dimension, even with arbitrary
finite-order multiplicity blocks. Thus the support fraction on
all nonzero linear weights is at least
\[
                (1-1/\ell)(1-q^{-r})
                 \ge(1-1/\ell)(1-q^{-r/2}).
\]
Give the trivial $Q_e$-isotypic summand its worst possible full
eigenspace. Adding the preceding bounds gives the first inequality
in (HOS5.7), and D-HOS-07 gives the second. This argument includes
$\ell=p$: neither semisimplicity of the natural action nor a
coprime averaging projector was used.
\end{proof}

\begin{lemma}[D-HOS-09: every compatible normalizer over $\mathbf F_5$]
\label{D-HOS-09}
Let $n\ge4$, and let $\rho$ be any nontrivial ordinary irreducible
representation of $G=\operatorname{Sp}_{2n}(5)$ of degree $a$.
Every nonscalar finite-order operator $T$ normalizing $\rho(G)$
satisfies
\[
                          \nu(T)\ge\frac{31}{250}a.           \tag{HOS5.11}
\]
This includes every compatible outer automorphism and both possible
central quotients of $G$.
\end{lemma}

\begin{proof}
First record three consequences of D-HOS-08. If $s$ is a natural
involution with minus-one eigenspace a nondegenerate plane, choose
$e,f$ in that plane. On the associated line radical, conjugation is
$(w,z)\mapsto(-w,z)$. Thus
\[
 \nu(\rho(s))\ge\frac{1-5^{-(n-1)}}4a
                          \ge\frac{31}{125}a.               \tag{HOS5.12}
\]
If $s$ is any noncentral natural involution, choose $e,f$ in its
smaller nonzero eigenspace. Both eigenspaces are nondegenerate and
even dimensional. Writing its eigenvalue on that plane as
$\lambda=\pm1$, the radical action is $w\mapsto\lambda s|_Ww$,
which moves the whole larger eigenspace of dimension $2m$ with
$m\ge\lceil n/2\rceil\ge2$. Therefore
\[
                         \nu(\rho(s))\ge\frac6{25}a.        \tag{HOS5.13}
\]

If $u=I+N\in G$ has $N^2=0$ and $\operatorname{rank}N=2$, its
image $R$ is isotropic and $N$ is skew-adjoint. To verify this
directly, apply preservation of $B$ by $I+jN$ for $j=1,2$:
subtracting twice the first identity from the second gives
$2B(Nx,Ny)=0$, and the first then gives skew-adjointness.
It follows that $\ker N=R^\perp$. Complete $R$ to a nondegenerate
four-dimensional symplectic subspace $E_0$; its orthogonal
complement lies in $\ker N$. Choose a paired $e,f$ in this
complement, which has dimension at least four. On its line radical
the induced action is $u|_W$, of order five and moving rank two.
Consequently
\[
                         \nu(\rho(u))\ge\frac8{25}a.        \tag{HOS5.14}
\]

We now pass to an arbitrary compatible normalizer. The ordinary
universal cover and automorphism inputs are exactly
\cite[Chapter~5, Table~1, p.~35]{BI:HHM} and
\cite[\S2.2.2, Table~2.1, author p.~14]{HES:LiXia2022}.
In this range the universal cover is $G$, and all automorphisms of
the simple quotient are inner-diagonal: the prime field has no
field automorphisms and type $C_n$, $n\ge4$, has no graph
automorphism. Thus the automorphism implemented by $T$ lifts to
conjugation by an $h\in\operatorname{GSp}(E)$; put
$\mu=\mu(h)\in\mathbf F_5^*$.
This assertion also holds when $\rho$ kills the central sign.
The kernel of a nontrivial irreducible representation is central,
and automorphisms of its image induce automorphisms of the simple
quotient. The natural lift and the given automorphism coincide:
their difference, if trivial on the simple quotient, would be a
central character of a perfect group. Such a character is trivial.
Hence
\[
                     T\rho(x)T^{-1}=\rho(hxh^{-1})
                            \quad(x\in G).                  \tag{HOS5.15}
\]

For $0\ne v\in E$ and $c\in\mathbf F_5^*$ use the actual
transvection (HOS5.4). Conjugation gives
\[
 h t_v(c)h^{-1}=t_{hv}(c/\mu),\qquad
 k=[h,t_v(c)]=t_{hv}(c/\mu)t_v(-c).                         \tag{HOS5.16}
\]
On the complex representation its commutator is $\rho(k)$.
For any chosen eigenvalue $\alpha$ of $T$,
\[
 \operatorname{rank}(\rho(k)-I)
 =\operatorname{rank}(T\rho(t_v(c))-\rho(t_v(c))T)
 \le2\operatorname{rank}(T-\alpha I).
                                                               \tag{HOS5.17}
\]
The equality discards invertible right factors. Choosing the
largest eigenspace bounds this by $2\nu(T)$.

Suppose first some $v$ has $b=B(v,hv)\ne0$. The plane
$P=\langle v,hv\rangle$ is nondegenerate, $k$ is the identity
on $P^\perp$, and in the indicated ordered basis its matrix is
\[
 k|_P=
 \begin{pmatrix}1&-cb\\-cb/\mu&1+c^2b^2/\mu\end{pmatrix},
 \qquad \det(k|_P)=1,\quad
 \operatorname{Tr}(k|_P)=2+c^2b^2/\mu.                     \tag{HOS5.18}
\]
If $\mu$ is a square, choose $c$ so that $c^2b^2/\mu=4$;
the trace is then one, and Cayley--Hamilton gives $k^3=-I$ on
$P$. If $\mu$ is a nonsquare, choose $c$ so that this ratio is
three; the trace is then zero, and $k^2=-I$ on $P$. These choices
exist because the nonzero squares in $\mathbf F_5$ are $1,4$
and the nonsquares are $2,3$. In either case a power of $k$ is
the rank-two involution in (HOS5.12). Therefore
\[
 2\nu(T)\ge\operatorname{rank}(\rho(k)-I)
          \ge\nu(\rho(k))
          \ge\frac{31}{125}a.                              \tag{HOS5.19}
\]

It remains that $B(v,hv)=0$ for every $v$. Polarizing this
identity gives $B(x,hy)=B(hx,y)$, so $h$ is self-adjoint for
$B$. Its adjoint as a similitude is $\mu h^{-1}$; consequently
$h^2=\mu I$.
If $\mu$ is a square, rescale $h$ to an involution $s\in G$.
Equation (HOS5.15) and Schur's lemma make $T$ a scalar multiple
of $\rho(s)$. If $s$ were central then $T$ would be scalar;
otherwise (HOS5.13) is stronger than (HOS5.11).

If $\mu$ is nonsquare, $h$ has no eigenline over $\mathbf F_5$.
For any $v\ne0$ the vectors $v,hv$ are independent and span a
totally isotropic plane. Equation (HOS5.16) now has the form
\[
 k=I+N,\qquad
 N=\frac c\mu hv\,B(hv,-)-c v\,B(v,-).
                                                               \tag{HOS5.20}
\]
The two pairing functionals and the two image vectors are
independent. Thus $\operatorname{rank}N=2$, while isotropy
gives $N^2=0$. Equations (HOS5.14) and (HOS5.17) give
$\nu(T)\ge4a/25>31a/250$. All cases have now been covered.
\end{proof}

\begin{theorem}[D-HOS-10: no higher $\mathbf F_5$ source]
\label{D-HOS-10}
In the genus-zero finite-group source problem for an imaginary-quadratic
Weil sixfold, no inducing representation whose image has an
irreducibly acting normal central product of nontrivial quasisimple
factors with simple quotients
\[
               \operatorname{PSp}_{2n_i}(5),\qquad n_i\ge4,
                                                               \tag{HOS5.21}
\]
can have period rank nine. The factors may have different ranks;
no transitivity on them is required. This conclusion concerns either
polarization determinant and remains valid if the inducing character
field is larger than the original imaginary quadratic field.
\end{theorem}

\begin{proof}
Let the inducing image be the finite faithful irreducible group
$J\le\operatorname{GL}(V)$ and put $e=\dim V$.
Restriction to its normal irreducible central product is a tensor
product of nontrivial ordinary irreducibles of the covers
$\operatorname{Sp}_{2n_i}(5)$. The elementary D-HOS-02 bound gives
\[
                   a_i\ge2\cdot5^{n_i-1}\ge250,
                      \qquad e=\prod_i a_i\ge250.           \tag{HOS5.22}
\]

First extend (HOS5.11) to all nonscalar elements of $J$. A
factor-preserving normalizer is a scalar times a tensor product
of compatible factor normalizers, by the full matrix-algebra
commutant. Its largest normalized eigenvalue multiplicity is at
most that of any one nonscalar factor. The factor operators can
be rescaled to finite order, since a power of their tensor
product is scalar. If an element permutes factors, consider any
nontrivial cycle of length $s\ge2$ on equal-dimensional factors
of degree $b\ge2$. By factorwise changes of basis its operator
is a cyclic rotation with one final diagonal twist. In the
tensor basis, the $b$ constant words are the only fixed words;
every other word orbit has at least two elements. A weighted
cyclic orbit contributes at most one vector to any one eigenspace.
Thus the largest eigenspace on that cycle has dimension at most
\[
                 b+\frac{b^s-b}{2}\le\frac34 b^s.
\]
Convolution with the remaining cycles preserves this upper bound.
In both cases every nonscalar element of $J$ has support at least
$31e/250$.

We next prove essential selfduality, including its source-relevant
nontrivial multiplier. The published factorization theorem of
Wonenburger, stated with its full field hypothesis in
\cite[Theorem~1, author p.~1]{HOS:Vinroot2004}, says that any
$x\in\operatorname{Sp}_{2n}(5)$ is a product of two
anti-symplectic involutions. If $x=h_1h_2$, then
$h_1xh_1=x^{-1}$. For a fixed anti-symplectic involution $h_0$,
the product $h_1h_0$ is symplectic, so $h_0xh_0$ is
symplectically conjugate to $x^{-1}$. Since $-1=2^2$ in
$\mathbf F_5$, the map $2h_0$ is itself symplectic and has the
same conjugation action. Therefore every $x$ is conjugate to
$x^{-1}$ in its symplectic group. All its ordinary irreducible
characters are real-valued, and their representations are
selfdual. No assertion about their Frobenius--Schur signs is needed.

Their tensor product has a unique nondegenerate invariant bilinear
form up to scalar, by irreducibility. Every normalizer in $J$
therefore preserves that line and has a bilinear similitude
character $\eta:J\to\mathbf C^*$. Let $L$ be the field of the
inducing character. The original character is induced from it,
so its values are sums of values in $L$; consequently the original
imaginary quadratic field $K$ is contained in $L$. If $\eta$
were trivial, the whole $J$-representation would be selfdual
and its character real-valued, contradicting $K\subset L$.
Thus $\eta$ is nontrivial.

We finish with the actual variation-of-Hodge-structure factorization.
The polarized Shapiro identification of D-BI-01 realizes the
six-dimensional multiplicity variation on a finite Hurwitz family
over the intermediate quotient curve $X$. If its genus is $h$
and its genuine branch inertias in $J$ are $g_1,\ldots,g_R$,
parabolic Euler characteristic gives
\[
        \sum_{j=1}^R\operatorname{rank}(g_j-I)
                        =6+2e(1-h).                         \tag{HOS5.23}
\]
Trivial inertia is omitted, so $R$ is at most this sum. For
$h\ge2$ the right side is negative because $e\ge250$. For
$h=1$ it is six, and the actual period map factors through a
finite Hurwitz cover of $\mathcal M_{1,R}$ of dimension at
most six. Rank nine therefore forces $h=0$, and then the
same actual factorization through $\mathcal M_{0,R}$ forces
$R\ge12$. These are bounds on the moduli space through which the
actual intermediate period map factors.

The genus-zero inertias generate $J$ and have product one.
Because $\eta$ is nontrivial, at least two have nontrivial
similitude. Their fixed subspaces are totally isotropic for the
nondegenerate bilinear form: for fixed vectors $v,w$,
$B(v,w)=B(gv,gw)=\eta(g)B(v,w)$. Each such inertia has moving
rank at least $\lceil e/2\rceil$. Every remaining genuine
inertia has moving rank at least $31e/250$: this follows from
the support bound if it is nonscalar, and from its full moving
rank $e$ if it is a nonidentity scalar. Applying (HOS5.23) to
two of the former and at least ten of the latter yields
\[
 2e+6\ge2\lceil e/2\rceil+10\frac{31e}{250}
          \ge\frac{56e}{25},\qquad e\le25.                 \tag{HOS5.24}
\]
This contradicts (HOS5.22) and proves the theorem.
\end{proof}

\paragraph{Certificate scope.}
The adjacent exact script checks the transvection trace and power
identities, all nonzero field parameters, the line-radical and
opposite-radical identities, the two self-adjoint multiplier types,
the constant inequalities, and the source-budget arithmetic. The
universal generation, carrier, normalizer, and source-factorization
proofs are the written arguments above. Finite matrix checks are
not substituted for them.
